\chapter{Olympiad geometry}\index{geometry!olympiad}

Olympiad geometry rewards \emph{configuration recognition}: the same
five or six lemmas --- cyclic quadrilaterals, the incenter/excenter
picture, power of a point, Ceva, homothety, inversion --- generate most
of what contests ask. This chapter teaches the tools, their equality
cases, and the escape hatch: when synthetic insight stalls, coordinates
and complex numbers are honest backup plans, not admissions of defeat.

% from-zero:begin
\begin{rulebox}[Before you start]
Chapter~2 (proofs) and school geometry: angles in a triangle, similar triangles, what a circle is. \emph{Recap:} the inscribed angle theorem says an angle at the circle is half the central angle on the same arc.
\end{rulebox}
\begin{wins}[After this chapter you can]
\begin{itemize}
\item find unknown angles systematically by angle chasing;
\item recognise cyclic quadrilaterals and use power of a point;
\item apply Ceva, Menelaus, homothety and inversion;
\item fall back on coordinates when a synthetic proof is not in sight.
\end{itemize}
\end{wins}
\begin{principle}[How to think about this]
Olympiad geometry rewards looking for the \emph{reason} an angle or a length is what it is. Draw a large, careful picture, mark every equal angle you can justify, and look for circles hiding in the figure. If you are stuck, you can always switch to coordinates, a slower but reliable backup.
\end{principle}
% from-zero:end

\section{Angle chasing systematically}\index{angle chasing}
School angle chasing uses unsigned angles and silently splits cases ---
acute vs obtuse, interior vs exterior. \emph{Directed
angles}\index{directed angle} work mod $180^\circ$:
$\measuredangle(\ell,m)$ is the turn from $\ell$ to $m$, so supplements
collapse into one statement (Figure~\ref{fig:directed}). The upgrade
costs one definition and pays back every case split.
\begin{figure}[ht]\centering
% Figure: directed-angle convention mod 180.
% Source for Figure~\ref{fig:directed}. TikZ only; regenerable by compiling main.tex.
\begin{tikzpicture}[
  lab/.style={font=\footnotesize\sffamily,text=slate,align=center},
]
\draw[ink,thick] (0,0) -- (2.2,0.8);
\draw[ink,thick] (0,0) -- (2.0,-0.7);
\draw[accent,thick,->] (0.9,0.03) arc (10:-19:0.95);
\node[lab] at (3.0,0.1) {$\measuredangle(\ell,m) \in \mathbb{R}/180^\circ\mathbb{Z}$:\\
$a$ and $a+180^\circ$ are the \emph{same} directed angle};
\node[lab,align=center]{no configuration cases ---\\supplements collapse};
\end{tikzpicture}

\caption{Directed angles are taken mod $180^\circ$. Notice what this
buys: the ``supplementary'' configuration stops existing.}
\label{fig:directed}
\end{figure}
The core rule: $\measuredangle(l,m) + \measuredangle(m,n) =
\measuredangle(l,n)$; and $\measuredangle(l,m)=0$ iff the lines
coincide.

\begin{example}[Reverse angle chase: prove a line is tangent]\index{tangent!angle chase}
Given $\triangle ABC$ with circumcircle $\omega$, show the line $\ell$
through $A$ with $\measuredangle(\ell, AB) = \measuredangle ACB$ is
tangent to $\omega$.

\textbf{Strategy.} Don't prove tangency head-on --- chase backwards
from the tangent criterion (a line is tangent iff its angle with the
chord equals the inscribed angle in the opposite arc).

\textbf{Solution.} Let $t$ be the tangent at $A$: then
$\measuredangle(t,AB)=\measuredangle ACB$ by the tangent--chord
theorem. But $\ell$ is defined by the \emph{same} angle on the same
side, so $\ell=t$: the constructed line \emph{is} the tangent. (Each
directed angle value determines a \emph{unique} ray through $A$ ---
matching the tangent's angle means matching the tangent's line.)
Reverse chasing converts ``prove tangent'' into ``match the tangent's
angle''.
\end{example}

\section{Cyclic quadrilaterals}\index{cyclic quadrilateral}
Half of olympiad geometry reduces to spotting (or building) a cyclic
quad: $ABCD$ cyclic iff $\measuredangle ACB=\measuredangle ADB$ (equal
angles on chord $AB$) iff $\angle A+\angle C=180^\circ$
(Figure~\ref{fig:cyclic}). Proving concyclicity is often the entire
problem in disguise.
\begin{figure}[ht]\centering
% Figure: cyclic quadrilateral -- equal angles on the same chord.
% Source for Figure~\ref{fig:cyclic}. TikZ only; regenerable by compiling main.tex.
\begin{tikzpicture}[
  lab/.style={font=\footnotesize\sffamily,text=slate,align=center},
  pt/.style={circle,fill=ink,inner sep=1.4pt},
]
\draw[ink] (0,0) circle (1.5);
\foreach \a/\n in {200/A,160/B,20/C,-30/D}
  \node[pt,label={[lab]180+\a:\n}] (\n) at (\a:1.5) {};
\draw[rulegray] (A) -- (C);
\draw[rulegray] (B) -- (C);
\draw[rulegray] (A) -- (D);
\draw[rulegray] (B) -- (D);
\draw[accent,thick] (A) -- (B);
\node[lab] at (0,-2.2) {$\angle ACB = \angle ADB$: angles standing\\
on chord $AB$ from the same arc agree};
\end{tikzpicture}

\caption{Equal angles on a chord force concyclicity. Notice the converse
is a construction tool: need two angles equal? Show the four points are
concyclic.}
\label{fig:cyclic}
\end{figure}

\begin{example}[Spot the cyclic quad first]
Acute $\triangle ABC$ has altitudes $BE, CF$ meeting at $H$. Prove
$A, F, H, E$ are concyclic and conclude $\angle FAH = \angle FEH$.
\par\smallskip
\textbf{Strategy.} Don't compute angles --- look for right angles
sharing a hypotenuse: two right angles on one segment force a circle
by Thales.
\par\smallskip
\begin{solution}
$\angle AFH = 90^\circ$ and $\angle AEH = 90^\circ$: both $F$ and $E$
see $AH$ under a right angle, so both lie on the circle with diameter
$AH$ --- $A,F,H,E$ concyclic. Equal angles on chord $FH$ now give
$\angle FAH = \angle FEH$ without any angle arithmetic. The move
repeats everywhere: altitudes, tangents and midpoints all manufacture
right angles, and right angles manufacture cyclic quads.
\end{solution}
\end{example}

\section{The incenter/excenter lemma}\index{incenter/excenter lemma}
\begin{lemma}
Let $I$, $I_A$ be the incenter and $A$-excenter of $\triangle ABC$, and
$M$ the midpoint of arc $BC$ not containing $A$ on the circumcircle.
Then $M$ is equidistant from $B,I,C,I_A$: $MB=MI=MC=MI_A$.
\end{lemma}
\begin{proof}[Proof sketch]
$\angle IBM$ splits as $\angle IBC + \angle CBM$; the first is $B/2$,
the second equals $\angle CAM = A/2$ since both subtend arc $CM$ \dots
the symmetric chase
gives $\angle IBM = \angle BIM$, so $MB=MI$; excenter goes the same way
with external bisectors (Figure~\ref{fig:incexc}).
\end{proof}
\begin{figure}[ht]\centering
% Figure: incenter/excenter lemma -- arc midpoint equidistant from B, I, C.
% Source for Figure~\ref{fig:incexc}. TikZ only; regenerable by compiling main.tex.
\begin{tikzpicture}[
  lab/.style={font=\footnotesize\sffamily,text=slate,align=center},
  pt/.style={circle,fill=ink,inner sep=1.4pt},
]
\draw[ink] (0,0) circle (1.6);
\node[pt,label={[lab]above:$A$}] (A) at (90:1.6) {};
\node[pt,label={[lab]left:$B$}] (B) at (200:1.6) {};
\node[pt,label={[lab]right:$C$}] (C) at (-20:1.6) {};
\node[pt,label={[accent]below:$M$}] (M) at (-90:1.6) {};
\node[pt,label={[lab]0:$I$}] (I) at (0.1,0.15) {};
\draw[rulegray] (A) -- (B) -- (C) -- cycle;
\draw[accent,thick] (M) -- (B);
\draw[accent,thick] (M) -- (I);
\draw[accent,thick] (M) -- (C);
\draw[rulegray,dashed] (A) -- (M);
\node[lab,align=center]{$MB = MI = MC$ (and $=MI_A$):\\
the arc midpoint is a stealth circle's center};
\end{tikzpicture}

\caption{The incenter/excenter lemma: arc midpoint $M$ is secretly the
center of a circle through $B,I,C,I_A$. Notice the pattern --- when
four named points turn out concyclic, a ``configuration'' is present.}
\label{fig:incexc}
\end{figure}

\section{Power of a point}\index{power of a point}
For point $P$ and circle of radius $r$ center $O$: every secant through
$P$ meets the circle in $A,B$ with $PA\cdot PB = OP^2-r^2$ (signed);
for $P$ outside this equals $PT^2$ for tangent $PT$; for $P$ inside,
$PA\cdot PB = r^2-OP^2$ (Figure~\ref{fig:power}). Two circles' equal-power
locus is the \emph{radical axis}\index{radical axis}, a line; three
circles give the \emph{radical center} where the three axes concur
(Figure~\ref{fig:radical}).
\begin{figure}[ht]\centering
% Figure: power of a point, three cases (interior / exterior / tangent).
% Source for Figure~\ref{fig:power}. TikZ only; regenerable by compiling main.tex.
\begin{tikzpicture}[
  lab/.style={font=\footnotesize\sffamily,text=slate,align=center},
  pt/.style={circle,fill=ink,inner sep=1.2pt},
]
% interior
\draw[ink] (-3.2,0) circle (1.0);
\node[pt,label={[lab]below:$Q$}] at (-3.2,0.1) {};
\draw[rulegray] (-4.1,-0.3) -- (-2.3,0.5);
\draw[rulegray] (-3.9,0.6) -- (-2.6,-0.7);
\node[lab] at (-3.2,-1.5) {$QA\cdot QB = r^2-OQ^2$};
% exterior
\draw[ink] (1.0,0) circle (1.0);
\node[pt,label={[lab]right:$P$}] at (2.9,0) {};
\draw[rulegray] (2.9,0) -- (0.3,0.55);
\draw[accent,thick] (2.9,0) -- (1.7,0.95);
\node[lab] at (1.6,-1.5) {$PA\cdot PB = PT^2 = OP^2-r^2$};
\end{tikzpicture}

\caption{Power of a point, inside and out. Notice the value depends
only on $OP$ --- the \emph{direction} of the secant is irrelevant, which
is what makes it useful.}
\label{fig:power}
\end{figure}
\begin{figure}[ht]\centering
% Figure: radical axis of two disjoint circles.
% Source for Figure~\ref{fig:radical}. TikZ only; regenerable by compiling main.tex.
\begin{tikzpicture}[
  lab/.style={font=\footnotesize\sffamily,text=slate,align=center},
]
\draw[ink] (-1.4,0) circle (0.9);
\draw[ink] (1.6,0) circle (1.2);
\draw[accent,thick] (0.35,-1.6) -- (0.35,1.6);
\node[lab] at (0.35,1.9) {radical axis};
\node[lab,align=center]{points with equal power w.r.t.\ both\\
--- construct it from two secants};
\end{tikzpicture}

\caption{The radical axis of two disjoint circles. Notice it exists
even when the circles don't meet --- ``common chord'' generalizes to
``equal-power locus''.}
\label{fig:radical}
\end{figure}

\begin{example}[Two-secant length computation]\index{power of a point!example}
Secants from $P$ cut a circle at $(A,B)$ and $(C,D)$ with $PA=3$,
$AB=5$, $PC=2$. Find $CD$.

\textbf{Strategy.} Pick the circle first, then the point --- the
products $PA\cdot PB$ and $PC\cdot PD$ are the same number.

\textbf{Solution.} $PA\cdot PB = 3\cdot 8 = 24 = PC\cdot PD = 2\cdot PD$
so $PD=12$ and $CD=10$. The answer is the arithmetic of one invariant.
\end{example}

\section{Ceva, Menelaus, extended sine law}\index{Ceva}\index{Menelaus}
Ceva: cevians $AD,BE,CF$ concur iff $\frac{BD}{DC}\cdot
\frac{CE}{EA}\cdot\frac{AF}{FB}=1$ (directed lengths; the product is
$+1$). Menelaus: a transversal cutting the three side-lines gives
product $-1$ (directed sign distinguishes the two theorems)
(Figure~\ref{fig:ceva}). The \emph{extended law of sines}: in any
triangle $a/\sin A = 2R$, connecting ratios to the circumradius.
\begin{figure}[ht]\centering
% Figure: Ceva and Menelaus side by side with directed-segment arrows.
% Source for Figure~\ref{fig:ceva}. TikZ only; regenerable by compiling main.tex.
\begin{tikzpicture}[
  lab/.style={font=\footnotesize\sffamily,text=slate,align=center},
  pt/.style={circle,fill=ink,inner sep=1.3pt},
]
% Ceva triangle
\node[pt,label={[lab]above:$A$}] (A) at (-3.2,1.4) {};
\node[pt,label={[lab]left:$B$}] (B) at (-4.4,-1.0) {};
\node[pt,label={[lab]right:$C$}] (C) at (-2.0,-1.0) {};
\draw[ink] (A) -- (B) -- (C) -- cycle;
\coordinate (D) at (-3.6,-1.0);
\coordinate (E) at (-2.6,0.2);
\coordinate (F) at (-3.9,0.3);
\draw[accent] (A) -- (D); \draw[accent] (B) -- (E); \draw[accent] (C) -- (F);
\foreach \p/\n in {D/D,E/E,F/F} \node[lab] at (\p) [below right=-2pt] {\tiny\n};
\node[lab] at (-3.2,-1.9) {Ceva: $\frac{BD}{DC}\frac{CE}{EA}\frac{AF}{FB}=1$};
% Menelaus triangle
\node[pt,label={[lab]above:$A'$}] (A2) at (1.6,1.4) {};
\node[pt,label={[lab]left:$B'$}] (B2) at (0.4,-1.0) {};
\node[pt,label={[lab]right:$C'$}] (C2) at (2.8,-1.0) {};
\draw[ink] (A2) -- (B2) -- (C2) -- cycle;
\draw[accent,dashed,thick] (0.6,-0.35) -- (3.6,-0.05);
\node[lab] at (1.6,-1.9) {Menelaus: a transversal gives $= -1$};
\end{tikzpicture}

\caption{Ceva and Menelaus side by side. Notice the directed-length
convention is what makes the two statements differ in sign rather than
in case lists.}
\label{fig:ceva}
\end{figure}
\begin{example}[Medians concur --- and a ratio chase]\index{medians!concurrent}
\textbf{Strategy.} For medians, every factor is $1$; for a nontrivial
case, write the product in directed form and identify each factor.

\textbf{Solution.} Medians: $BD=DC$ etc., product $1$, so the three
medians concur. Ratio chase: if $D,E,F$ divide sides in ratios
$x,y,z$ cyclically, concurrence iff $xyz=1$ --- e.g.\ cevians dividing
sides in ratios $1/2, 3, 2/3$ concur since
$\tfrac12\cdot3\cdot\tfrac23=1$.
\end{example}

\section{Homothety and spiral similarity}\index{homothety}
A homothety (dilation) at center $Z$ ratio $k$ sends every figure to a
similar one; a \emph{spiral similarity} composes it with a rotation ---
and is the unique map taking one segment $AB$ to another $A'B'$
directly. Circles tangent internally are homothety images of each other.

\begin{example}[The nine-point circle exists]\index{nine-point circle}
\textbf{Strategy.} Don't construct nine points --- find one dilation:
the homothety at the orthocenter $H$ with ratio $\tfrac12$ should send
the circumcircle onto the claimed circle.

\textbf{Solution.} The homothety $h(H,\tfrac12)$ sends each vertex to
the midpoint of its $H$-segment and the circumcircle to a circle
$\Gamma$ through the three feet of the altitudes and the three
midpoints of the sides' Euler segments; the homothety $h(G,-\tfrac12)$
at the centroid sends the circumcircle to the circle through the side
midpoints --- and both images coincide since $N$ is the midpoint of
$OH$ (Figure~\ref{fig:ninept}). Existence is a corollary of finding the
dilation; the nine points ride for free.
\end{example}
\begin{figure}[ht]\centering
% Figure: nine-point circle as homothety image of the circumcircle at H, ratio 1/2.
% Source for Figure~\ref{fig:ninept}. TikZ only; regenerable by compiling main.tex.
\begin{tikzpicture}[
  lab/.style={font=\footnotesize\sffamily,text=slate,align=center},
  pt/.style={circle,fill=ink,inner sep=1.3pt},
]
\draw[ink] (0,0) circle (1.8);
\draw[accent,dashed] (-0.3,-0.2) circle (0.9);
\node[pt,label={[lab]left:$O$}] (O) at (0,0) {};
\node[pt,label={[lab]right:$H$}] (H) at (0.9,-0.6) {};
\node[pt,label={[accent]below:$N$}] (N) at (-0.3,-0.2) {};
\draw[rulegray] (O) -- (H);
\node[lab] at (0,-2.5) {homothety at $H$, ratio $\tfrac12$:\\
circumcircle $\to$ nine-point circle; $N$ is midpoint of $OH$};
\end{tikzpicture}

\caption{The nine-point circle as homothety image. Notice the proof is
one map, not nine verifications --- the dilation \emph{is} the theorem.}
\label{fig:ninept}
\end{figure}

\section{Inversion in a circle}\index{inversion}
Inversion in circle $(O,r)$ sends $P\mapsto P'$ on ray $OP$ with
$OP\cdot OP' = r^2$. It swaps inside/outside, reverses orientation, and
converts between ``lines'' and ``circles through $O$''
(Figure~\ref{fig:invert}): lines $\not\ni O \leftrightarrow$ circles
$\ni O$; circles $\not\ni O \leftrightarrow$ circles $\not\ni O$.
Choose the center where the problem's concurrency or tangency lives.
\begin{figure}[ht]\centering
% Figure: inversion image table (line<->circle through O etc.).
% Source for Figure~\ref{fig:invert}. TikZ only; regenerable by compiling main.tex.
\begin{tikzpicture}[
  lab/.style={font=\footnotesize\sffamily,text=slate,align=center},
]
\draw[ink,dashed] (0,0) circle (0.8);
\node[circle,fill=ink,inner sep=1.2pt,label={[lab]below:$O$}] (O) at (0,0) {};
\draw[rulegray,thick] (-1.9,1.1) -- (1.9,1.1);
\draw[accent,thick] (0,0.55) circle (0.28);
\node[lab] at (-1.9,1.4) {line};
\node[lab] at (1.5,0.55) {circle thru $O$};
\node[lab,align=center]{line $\not\ni O \;\leftrightarrow\;$ circle $\ni O$;\\
circle $\not\ni O \;\leftrightarrow\;$ circle $\not\ni O$};
\end{tikzpicture}

\caption{Inversion's image table. Notice what it trades: your mess of
circles becomes lines --- and tangency (a hard thing to use) becomes
parallelism or collinearity (easy).}
\label{fig:invert}
\end{figure}
Worked transformation: to prove four circles share a common point,
invert about one circle's center; the four become lines or circles
through $O$ whose concurrency is a collinearity check. This chapter
teaches recognition and one deployment; the deeper problem set lives in
the Notes.

\section{Analytic backup plans}\index{analytic geometry}
Coordinates, areas, trigonometric form, and Ptolemy's inequality
($AC\cdot BD \le AB\cdot CD + BC\cdot AD$ for any quadrilateral,
equality iff cyclic) are the backup plan --- honest, mechanical, and
sometimes strictly cheaper than synthetic insight. Fix scale and
orientation first (e.g.\ unit circle at origin); choose analytic when
the problem is asymmetric or the synthetic handle is invisible.

\section{Summary}
\begin{center}\small
\begin{tabular}{@{}ll@{}}
\toprule
\textbf{Tool} & \textbf{What it buys} \\
\midrule
Directed angles & no case splits \\
Cyclic quads & equal angles $\leftrightarrow$ concyclic \\
Incenter/excenter & hidden circle centers \\
Power of a point & $PA\cdot PB$ invariant \\
Ceva/Menelaus & concurrency/collinearity $=\pm1$ \\
Homothety & one map $=$ existence proof \\
Inversion & circles $\leftrightarrow$ lines \\
Coordinates/complex & the backup plan \\
\bottomrule
\end{tabular}
\end{center}

\section{Code: two numerical verifications}\index{power\_check}\index{invert.py}
\texttt{power\_check.py} places coordinates and verifies
$PA\cdot PB = PT^2 = OP^2-r^2$ plus the interior variant.
\texttt{invert.py} inverts a sampled line and checks its image lies on
the predicted circle through $O$.
\begin{checkpoint}
Predict the common value for $P=(2,1)$ on the unit circle, and which
circle the line $x=1$ maps to under $r=2$ inversion. Then run both.
\end{checkpoint}
\begin{lstlisting}
secant: PA*PB = 4.000000
tangent: PT^2  = 4.000000
agree: True
interior: QA*QB = 0.870000 vs r^2-|OQ|^2 = 0.870000
line x=1 -> 13 image points
max distance from predicted circle (center (2,0), r=2): 4.44e-16
inv(inv(P))=P check: (2.999..., 1.0) vs (3.0, 1.0)
|OP|*|OP'| = 4.000 = r^2 = 4.0
\end{lstlisting}
\texttt{test\_15\_geom.sh} re-runs and greps the agreement lines.

\section{Common mistakes}\index{geometry!mistakes}
\begin{itemize}
\item \emph{Unsigned angles landing on the supplement} --- the
EGMO-errata mistake; directed angles exist precisely to prevent it.
\item \emph{Ceva/Menelaus on unsigned ratios} --- the sign is the whole
point; undirected products can't tell the theorems apart.
\item \emph{Assuming concyclicity} --- three points always define a
circle, four almost never do; prove it before using it.
\item \emph{Wrong inversion center} --- the circle of inversion is a
choice; put it at the concurrency or tangency point, and remember
orientation flips.
\item \emph{Unfixed coordinates} --- choosing axes before fixing scale
and orientation doubles the algebra.
\end{itemize}

\section*{Exercises}
\subsection*{Inline}
\begin{exercise}[easy]\label{exr:g1}
Show two parallel chords of a circle intercept equal arcs.
\end{exercise}
\begin{solution}
Inscribed angles on equal arcs; directed angles give it in one line.
\end{solution}
\begin{exercise}[easy]\label{exr:g2}
Secant: $PA=4$, $AB=6$; find $PT$ for a tangent from $P$.
\end{exercise}
\begin{solution}
$PT^2 = 4\cdot10 = 40$, $PT=2\sqrt{10}$.
\end{solution}
\begin{exercise}[medium]\label{exr:g3}
Prove the tangent--chord theorem used in the first worked example.
\end{exercise}
\begin{solution}
Draw the diameter through the contact point; the tangent is
perpendicular; inscribed-angle bookkeeping finishes.
\end{solution}
\begin{exercise}[medium]\label{exr:g4}
Show the three radical axes of three circles (taken pairwise) concur at
the radical center --- when they are not parallel.
\end{exercise}
\begin{solution}
A point on two axes has equal power w.r.t.\ all three circles, so it
lies on the third axis too.
\end{solution}
\begin{exercise}[medium]\label{exr:g5}
Ceva: points $D,E,F$ on $BC,CA,AB$ with $BD/DC=2$, $CE/EA=1/2$,
$AF/FB=1$. Concur?
\end{exercise}
\begin{solution}
$2\cdot\tfrac12\cdot1 = 1$: yes.
\end{solution}
\begin{exercise}[medium]\label{exr:g6}
Inversion: what is the image of a circle through $O$ under inversion
about $(O,r)$?
\end{exercise}
\begin{solution}
A line not through $O$ --- the ``missing point'' $O$ on the circle
opens it into a line.
\end{solution}
\subsection*{Additional}
\begin{exercise}[easy]\label{exr:ga1}
Two secants: $PA\cdot PB = PC\cdot PD$ implies $A,B,C,D$ concyclic ---
state and prove the converse.
\end{exercise}
\begin{solution}
Equal powers put $P$ on both circles' radical axis --- the converse of
power of a point gives the circle through $A,B,C,D$.
\end{solution}
\begin{exercise}[easy]\label{exr:ga2}
Compute the power of the circle's own center $O$.
\end{exercise}
\begin{solution}
$-r^2$ (signed) --- the minimum possible power.
\end{solution}
\begin{exercise}[easy]\label{exr:ga3}
Angles on the same arc are equal --- give the directed-angle one-liner.
\end{exercise}
\begin{solution}
$\measuredangle ACB=\measuredangle ADB$ iff $A,B,C,D$ concyclic; the
proof is the inscribed-angle theorem with signs.
\end{solution}
\begin{exercise}[easy]\label{exr:ga4}
A circle through $O$ inverts to a line --- what happens to $O$ itself?
\end{exercise}
\begin{solution}
$O$ has no image ($OP\cdot OP'=r^2$ is degenerate at $0$); the line is
the circle ``minus $O$'' closed at infinity.
\end{solution}
\begin{exercise}[medium]\label{exr:ga5}
Use the incenter/excenter lemma to find the center of the circle through
$B,I,C,I_A$ and its radius in terms of $M$'s position.
\end{exercise}
\begin{solution}
Center $M$; radius $MB$ --- computable from $2R\sin(A/2)$.
\end{solution}
\begin{exercise}[medium]\label{exr:ga6}
Ceva form of concurrency for angle bisectors: verify the incenter via
the bisector theorem.
\end{exercise}
\begin{solution}
By the angle-bisector theorem, $BD/DC=c/b$, and likewise for the other
two bisectors; the product $\frac cb\cdot\frac ac\cdot\frac ba = 1$
satisfies Ceva.
\end{solution}
\begin{exercise}[medium]\label{exr:ga7}
Show a spiral similarity taking $AB\mapsto CD$ has center at the
intersection of circles through $(A,C)$ and $(B,D)$ --- the
``pivot theorem'' idea.
\end{exercise}
\begin{solution}
The center sees $AC$ and $BD$ under the same rotation angle; the two
circles' second intersection is it.
\end{solution}
\begin{exercise}[medium]\label{exr:ga8}
Ptolemy: prove $AC\cdot BD = AB\cdot CD + BC\cdot AD$ for cyclic $ABCD$
by the inscribed-angle trick on the diagonal.
\end{exercise}
\begin{solution}
Take $E$ on $AC$ with $\angle ABE=\angle DBC$; two similar triangles
give the two products.
\end{solution}
\begin{exercise}[medium]\label{exr:ga9}
Implement-or-derive: extend \texttt{power\_check.py} to the interior
case for a \emph{random} chord through $Q$, 100 trials.
\end{exercise}
\begin{solution}
Loop over random directions; every trial's $QA\cdot QB$ lands on
$r^2-|OQ|^2$ to machine precision.
\end{solution}
\begin{exercise}[hard]\label{exr:ga10}
Prove the radical axis is a line, not just a locus.
\end{exercise}
\begin{solution}
$|PO_1|^2-r_1^2 = |PO_2|^2-r_2^2$ expands to a linear equation in
$P$'s coordinates --- the quadratic parts cancel.
\end{solution}
\begin{exercise}[hard]\label{exr:ga11}
Inversion converts the statement ``two circles are orthogonal'' into
what on their images?
\end{exercise}
\begin{solution}
Angle-preserving (conformal) maps preserve orthogonality; a circle
orthogonal to the inversion circle maps to itself.
\end{solution}
\begin{exercise}[hard]\label{exr:ga12}
Derive the nine-point circle radius in terms of $R$ and locate its
center on the Euler line.
\end{exercise}
\begin{solution}
Radius $R/2$ (the $\tfrac12$ homothety); center $N$ at the midpoint of
$OH$, between $G$ and $H$ on the Euler line $O$--$G$--$N$--$H$.
\end{solution}
\begin{exercise}[easy]\label{exr:ga13}
In cyclic $ABCD$, $\angle ABD = 40^\circ$. Find $\angle ACD$.
\end{exercise}
\begin{solution}
$40^\circ$ --- equal angles on chord $AD$.
\end{solution}
\begin{exercise}[medium]\label{exr:ga14}
Point $P$ inside a circle has two chords through it: $AB$ with
$PA=4,PB=9$ and $CD$ with $PC=3$. Find $PD$, and explain which theorem
is doing the work.
\end{exercise}
\begin{solution}
Intersecting-chords case of power of a point: $PA\cdot PB = PC\cdot PD$
gives $PD = 36/3 = 12$.
\end{solution}
\begin{exercise}[hard]\label{exr:ga15}
Show the midpoints of the three sides of a triangle and the feet of
the three altitudes are concyclic (the cheap half of the nine-point
circle) using only the cyclic-quad criterion twice.
\end{exercise}
\begin{solution}
Let $M_a,M_b,M_c$ be side midpoints and $D,E,F$ the feet. The midpoint
quadrilateral $M_aM_bM_c$ is homothetic to $ABC$ at ratio $1/2$ about
the centroid; for the feet: $M_bM_c \parallel BC$ and
$\angle BEC = \angle BFC = 90^\circ$ puts $E,F$ on the circle of
diameter $BC$ centered at $M_a$; combine with the trapezoid
$M_bM_cEF$ being isosceles-cyclic to close.
\end{solution}

\section*{Notes}
The configuration-first pedagogy and directed-angles-early order follow
EGMO's table of contents and its errata~\cite{egmo}; the
incenter/excenter lemma, the pivot theorem and the inversion image table
follow Chen's geometry handouts~\cite{chen-geom}. Zhao's
cyclic-quadrilateral, power-of-a-point and ``Three Lemmas'' handouts
carry the full inversion problem set routed here by the merge
decision~\cite{zhao}. Ceva/Menelaus in directed form and the analytic
backups follow Engel~\cite{engel}. Contest uses: 2024 AIME I P10, P5,
P8, P12; 2023 AIME I P5, P8; 2023 AIME II P9, P12; 2022 AIME I P8; 2021
AIME I P11, P13; 2020 AIME I P1, P15; 2018 AIME I P4, P13, P15; 2015
AIME I P11; 2025 AIME II P1, P5/P6; 2024 AMC 12A P19; USAJMO 2024 P1,
P6; 2023 USAJMO P2, P6.

\section*{Further reading}
\begin{itemize}
\item E.~Chen, \emph{Euclidean Geometry in Mathematical Olympiads}
(EGMO) --- the reference this chapter's order follows; Ch.~1--3 cover
everything here except inversion.
\item E.~Chen's geometry handouts ---
\url{https://web.evanchen.cc/handouts/}
--- free PDFs: incenter/excenter
lemma, pivot theorem, directed angles.
\item Y.~Zhao's olympiad geometry handouts ---
\url{https://yufeizhao.com/olympiad/}
--- cyclic quadrilaterals, power
of a point, and ``Three Lemmas in Geometry''.
\item A.~Engel, \emph{Problem-Solving Strategies}, geometry chapters
--- Ceva, Menelaus and the analytic backup plans.
\item R.~Hammack, \emph{Book of Proof} --- for writing
geometric arguments as clean proofs:

\url{https://richardhammack.github.io/BookOfProof/}
\item The AoPS community's geometry contest collections ---
every cited AIME/USAJMO item has community solution threads there:

\url{https://artofproblemsolving.com}
\end{itemize}

\begin{rulebox}[Watch instead]
\begin{itemize}
\item \textbf{Michael Greenberg --- ``Olympiad Geometry Problem \#93:
Two Isosceles Triangles''}, 12:25,
\url{https://www.youtube.com/watch?v=Hz8kyLKP0_k}
--- a spiral
similarity manufactures a cyclic quad exactly as the
cyclic-quadrilaterals and homothety sections teach.
(Verified 2026-10-11.)
\item \textbf{Sohil Rathi --- ``Circular Geometry, Power of a point,
Cyclic Quadrilaterals''}, 1:33:35,
\url{https://www.youtube.com/watch?v=NsQbhYfGh1Q}
--- an AMC-level class
lecture covering power of a point and cyclics live on problems.
(Verified 2026-10-11.)
\end{itemize}
\end{rulebox}

\chapternotes
